\documentclass[11pt]{article}
\usepackage[a4paper,margin=22mm]{geometry}
\usepackage{amsmath,amssymb,amsthm,booktabs,array}
\newtheorem{theorem}{Theorem}
\newtheorem{lemma}[theorem]{Lemma}
\newcommand{\HS}{\mathrm{HS}}
\newcommand{\Aut}{\operatorname{Aut}}
\newcommand{\Cor}{\operatorname{Cor}}
\newcommand{\Sym}{\operatorname{Sym}}
\title{Two explicit Higman--Sims incidence constructions}
\author{}\date{}
\begin{document}\maketitle
\vspace{-12mm}
\noindent
This component verifies the two constructions in Section~4 of Leemans
\cite{L}.  It uses the literal Higman--Sims group and the accepted finite
coset catalogue, and provides explicit maps between that catalogue and
the set constructions.  The historical Higman model and its identification
with $\HS$ are attributed in \cite{L} to Higman and Sims.  We do not replace
that attribution by a claim of a new construction.

Our retained group input is a faithful permutation model $G\cong\HS$ of
order $44\,352\,000$, together with its complete automorphism action of
order $88\,704\,000$.  The two previously recognized coset families have
parabolic orders
\[
 (2880,252000,40320),\qquad (2880,252000,252000).                 \tag{1}
\]
Their catalogue identifiers are 35 and 36.  These identifiers refer to
the acquired catalogue, not to numbering in the historical supplement.
We retain the classical identification of the unique Moore graph of
valency seven on fifty vertices as the Hoffman--Singleton graph, and its
automorphism group $U_3(5):2$, as used in \cite[Section~4]{L}.
The exact finite permutation and graph algorithms used in the certificate
are stated with their software inputs; group orders alone are never used
to identify an unnamed abstract group.

\begin{theorem}\label{thm:main}
There are sets $\mathcal P,\mathcal C,\mathcal Q$, of respective sizes
$176,1100,176$, in the retained literal model, with conics
$C\in\mathcal C$ eight-subsets and quadrics $Q\in\mathcal Q$
fifty-subsets of $\mathcal P$, satisfying all six Higman properties
printed in \cite[Section~4.1]{L}.  Put
$\mathcal L=\binom{\mathcal P}{2}$.
\begin{enumerate}
\item The geometry $\Gamma_C$ on $\mathcal P,\mathcal L,\mathcal C$,
with incidence given by containment, is isomorphic to family 35.
Its point residue is the Hoffman--Singleton graph, its conic residue is
the vertex--edge geometry of $K_8$, and its line residue is a digon.
It has $\Aut(\Gamma_C)=\Cor(\Gamma_C)=G$.
\item In $\Gamma_Q$ on $\mathcal P,\mathcal L,\mathcal Q$,
point incidence is containment.  A line is incident with a quadric
precisely when the quadric contains both conics through its two points.
This geometry is isomorphic to family 36.  Both point and quadric
residues are Hoffman--Singleton graphs, and its line residue is a digon.
Its type-preserving group is $G$, and $\Cor(\Gamma_Q)\cong\Aut(\HS)$.
\end{enumerate}
Both geometries have $61\,600$ chambers and Borel subgroup $S_6$.
Their line-residue kernels are $S_6$, and their line stabilizers are
split extensions $S_6:2^2$.  Their point stabilizers, and the quadric
stabilizers of $\Gamma_Q$, act faithfully as the full $U_3(5):2$
on their Hoffman--Singleton residues.  The conic stabilizer of
$\Gamma_C$ acts faithfully as $S_8$ on its eight points.
\end{theorem}

\section*{Literal Higman data and the sixth property}
Take the right-coset action of $G$ on a specified subgroup of order
$252000$.  The saved generator images give a homomorphism with trivial
kernel on the resulting 176 labelled points.  A saved element of the
complete ambient automorphism action conjugates the second family's
point parabolic to the first family's point parabolic.  This actual
element, all three transported subgroups, and all point-generator images
are checked.  Thus the two coordinate systems agree by an explicit map.

The conics and quadrics are the full orbits of the specified base
eight-set and fifty-set.  Every generator preserves the complete lists
and their incidence relations.  Each of the four sets of objects
(points, lines, conics and quadrics) is a single orbit.  Literal subset
calculations give fifty quadrics through each point and eight quadrics
through each conic.  For each of all $15400$ point pairs there are
exactly two conics $C_1,C_2$, and
\[
 \{Q:p,p'\in Q\}=\{Q:C_1\subseteq Q\}\cup\{Q:C_2\subseteq Q\},
 \qquad |\{Q:C_1\cup C_2\subseteq Q\}|=2.                  \tag{2}
\]
The two quadrics in the last set give a bijection from point pairs to
quadric pairs.  Conversely, every quadric pair contains exactly two
conics; its point intersection is their union, and their intersection
has size two.  These checks establish properties (i)--(v), including
both universal alternatives in (iv) and (v).

For property (vi), fix the retained conic $S$.  For each of its eight
points $p$ and each of its eight incident quadrics $Q$, test both
conditions in the source: every conic meeting $S$ in two points including
$p$ is contained in $Q$, and every conic contained in two quadrics
through $S$, one of which is $Q$, contains $p$.  The complete $8$ by $8$
truth table has exactly one simultaneous solution in each row and
column.  With the retained labels the matching is
\[
 (1,6),(26,37),(36,23),(61,34),(165,9),(167,17),(170,30),(174,41).
\]
The first coordinate is a point label and the second a quadric index.
Since the conditions use only incidence and $G$ is transitive on conics,
this proves (vi) for every conic.  The actual actions on points and on
quadrics are doubly transitive; the point stabilizer has quadric orbits
of lengths $50$ and $126$.  This also proves the source's transitivity
assertions on incident and nonincident point--quadric pairs.

\section*{Coset maps, flags, and residues}
For either family write $H_1,H_2,H_3$ in line, point, plane order.
The retained base objects $b_i$ have exactly these stabilizers in $G$.
For every coset the certificate records the map
\[
             H_i g\longmapsto b_i^g.                       \tag{3}
\]
It verifies each image literally and verifies the resulting six full
bijections.  Stabilizer equality makes (3) well-defined and equivariant.
Two right cosets are incident when they intersect.  The neighbours of
$H_i$ of type $j$ are therefore the orbit of $H_j$ under $H_i$.
For all twelve ordered pairs of distinct types in the two families,
the actual orbit is compared to the entire geometric neighbour set of
$b_i$.  Equivariance transports these equalities to all objects.  This
proves both incidence isomorphisms, without inferring incidence from
matching degrees alone.

Every line has two points and two incident planes, and all four
point--plane flags in its residue occur.  Every incident point--plane
pair extends through exactly seven lines.  Thus every partial flag
extends to a chamber, and both geometries have $15400\cdot4=61600$
chambers.  Their base chamber stabilizers have order $720$ and have
explicit bijective generator maps onto $\Sym(6)$.  Their $G$-orbits
therefore contain all chambers.  This is actual flag transitivity.
The previously recognized residual connectivity, local primitivity,
weak primitivity and intersection property transfer across (3).

For $\Gamma_C$, the base point residue has fifty conics as graph
vertices and 175 lines as graph edges.  All edges have distinct
two-element endpoint sets.  Breadth-first calculations at every vertex
give valency seven, diameter two and girth five.  The same calculation
holds for the point and quadric residues of $\Gamma_Q$.  The retained
uniqueness statement identifies these graphs as Hoffman--Singleton
graphs.  The conic residue of $\Gamma_C$ is literally its eight points
and all their 28 pairs, hence $K_8$.  Line residues are the complete
bipartite graph on two points and two planes.  Object transitivity
transports the base-residue conclusions to all residues.

The faithful actions of the three order-$252000$ parabolics on their
fifty-vertex residues have image order $252000$.  They therefore equal
the full known automorphism group of the Hoffman--Singleton graph.
The conic stabilizer's action equals $\Sym(8)$ and has trivial kernel.
For each line stabilizer, the kernel on the four residue objects is
exactly the already identified Borel $S_6$.  Its image is elementary
abelian of order four.  The certificate supplies two commuting
involutions generating an order-four subgroup disjoint from this kernel;
together they generate the full line stabilizer.  This proves the
splitting asserted by the colon notation, in addition to the kernel
and quotient orders.

\section*{The full automorphism and correlation groups}
For $\Gamma_C$, restrict an automorphism to the bipartite point--conic
incidence graph, with its two vertex classes distinguished.  This
restriction is faithful: a line is determined by its two points.
Every graph automorphism extends uniquely to lines, and the certificate
checks all returned generators on every line--conic incidence.
The complete graph calculation, on 1276 vertices, gives order
$44\,352\,000$.  The explicit $G$-generators lie in that group and
generate a faithful image of the same order, so equality holds.  Since
the three geometry types have distinct cardinalities, every correlation
preserves types.  Hence $\Cor(\Gamma_C)=\Aut(\Gamma_C)=G$.

For $\Gamma_Q$, use the point--quadric incidence graph, adjoining one
marker for each type.  The two markers form one separate colour class;
each is joined to all 176 objects of its type.  The 354-vertex graph
therefore allows precisely type-preserving or type-interchanging
permutations, not arbitrary permutations of incidence types.
The complete group has order $88\,704\,000$.  Every returned generator
is checked on all graph adjacencies.  It has a unique lift to lines:
if types are fixed, use the image point pair; if types are interchanged,
use the inverse of the point-pair/quadric-pair bijection in (2).
In both cases all $15400$ resulting line incidences are checked, so these
lifts are correlations of the full geometry.  Conversely every
correlation restricts to this graph, and the restriction is faithful.

Let $D$ be the resulting group and let $E$ be the explicit image of $G$.
The exact generator map $G\to E$ is bijective.  The kernel of $D$ on
the two markers equals $E$, its marker image has order two, and
$C_D(E)=1$.  Thus conjugation embeds $D$ in $\Aut(E)$; its order equals
the retained full order of $\Aut(\HS)$, so this embedding is surjective.
Consequently the type-preserving group is $G$ and the correlation group
is $\Aut(\HS)$.  A saved generator interchanges the types and provides
the duality claimed in \cite[Section~4.2]{L}.

\section*{Certificate boundary}
The finite files retain the complete object lists, both generator
actions, six coset maps, twelve neighbour comparisons, the sixth-axiom
matching, all small-graph generators and their line lifts, and explicit
Borel and complement maps.  Completeness of the graph-automorphism
calculations is an acquired nauty premise; the returned generators are
additionally checked literally.  The GAP finite group and homomorphism
algorithms and the retained identification of the ambient $\HS$ model
are explicit dependencies.  The original failed test of surjectivity
onto the declared range of an action homomorphism is preserved.  The
successor uses its trivial kernel, the correct faithfulness assertion.

This proves the two construction and group descriptions in Section~4
of the fixed six-page article.  It does not assert completion of the
article's enumeration, Table~1, or historical supplement labels.
Those obligations belong to their separately scoped components.

\begin{thebibliography}{9}
\bibitem{L} D. Leemans, \emph{The residually weakly primitive geometries
of HS}, fixed six-page source, pp.~231--236, Section~4, pp.~233--235.
\bibitem{H} G. Higman, \emph{On the simple group of D. G. Higman and
G. C. Sims}, Illinois J. Math. \textbf{13} (1969), 74--80.
The historical attribution is retained as stated in \cite{L}.
\end{thebibliography}
\end{document}
